\begin{abstract}
We train a reinforcement learning agent that decides, token by token, how long a small draft model should keep drafting before a large target model verifies its guesses. This speeds up LLM text generation while keeping the output identical to the large model's.

Speculative decoding pairs a small, fast draft model with a large target model: the draft proposes several tokens and the target checks them all in one forward pass. The number of tokens drafted per round decides the speedup. Too few and the target is called too often; too many and most drafts are thrown away. Current systems use a fixed draft length or a hand-set confidence threshold.

We model the drafting decision as a Markov Decision Process (MDP) whose state is built from the draft model's confidence signals and whose reward is tokens gained minus the cost of a target call. We solve it with four algorithms that span the course syllabus: Value Iteration (dynamic programming), Q-learning (temporal difference), Deep Q-Network (function approximation and deep RL) and REINFORCE (policy gradient). We train them in a fast simulator built from logged model outputs and evaluate them in a live speculative-decoding loop on two open model pairs (Qwen2.5 0.5B$\to$1.5B and SmolLM2 135M$\to$1.7B) on consumer GPUs, against fixed-length and threshold baselines. All four agents beat every baseline in real wall-clock time. The best, Value Iteration, generates
1.51$\times$ (Qwen) and 1.75$\times$ (SmolLM2) faster than the target model alone, up to 2.03$\times$ on code,
against 1.42$\times$ and 1.65$\times$ for HuggingFace's confidence-threshold rule. The agents rediscover the
threshold-style stopping rule that theory predicts is optimal, and learn on their own to draft longer on code
and maths than on chat.
\end{abstract}
